\section{LP Loading and Recovery Lower Bounds}
\label{sec:lp-hardness}

This section gives the LP thread of the paper.  Its endpoint is a
linear-size, bounded-degree standard-form LP whose reduced primal Newton
matrix is a scalar identity at a public infeasible start, whose reduced
log-barrier Hessian is a scalar identity along the entire central path, and
whose original-primal Newton-direction state nevertheless needs a linear
number of raw coefficient queries.  The obstruction is not the scalar solve:
it is construction of the hidden affine translation and recovery in the
original coordinates.

All state, access, residual, and success conventions are those of
\cref{sec:prelim-quantum-access,sec:prelim-output-contracts}.  Throughout
this section \(N\geq2\).  In particular, a residual without another
qualifier is the global RHS-relative residual
\(\norm{Ax-b}_2/\norm b_2\), and an output state is unconditional.  The same
results hold for a heralded branch of input-independent constant success
probability when repetitions are charged.  A reduced-coordinate state is not
an original-coordinate output, and an unobserved ``good branch'' is not an
output contract.  We count all raw queries used for setup, preconditioning,
right-hand-side formation, solution, and recovery.

The underlying amplification motifs are established.  Parity and the
polynomial method are classical query-complexity tools
\cite{beals2001-quantum-lower-bounds-by-polynomials}; endpoint padding and
biased clocks occur in HHL and Feynman--Kitaev constructions
\cite{harrow2009-quantum-algorithm-linear-systems-equations,caha2018-feynman-kitaev-clock,bausch2018-analysis-limitations-modified-circuit};
and sparse QLS states carrying an endpoint bit occur in Wang--Zhang and Mori
et al.\ \cite{wang2024-tight-quantum-depth-lower-bound,mori2026-sparsity-dependent-complexity-lower-bound}.
Apers--Gribling Theorem~8.4 already gives a linear exponent for sparse-LP
\emph{value} estimation in balanced dimensions
\cite{apers2026-quantum-speedups-linear-programming-interior}.  None of these
results is a constant-condition QLS lower bound of the kind sometimes
informally inferred from the statements below.  Our contribution is the
LP-native conjunction of regularity, reduced geometry, original-coordinate
output contracts, sharp residual thresholds, and complete access accounting.

\subsection{Parity amplification and feasible-state hardness}
\label{sec:lp-foundation}

Let \(\sigma_1,\ldots,\sigma_N\in\{\pm1\}\), and put
\[
 p_0=1,\qquad p_i=\prod_{j=1}^i\sigma_j.
\]
Represent a signed scalar at node \(j\) by a nonnegative pair
\((u_j,v_j)\), with \(d_j=u_j-v_j\) and \(q_j=u_j+v_j\).  A chain fixes
\begin{equation}
 d_0=1,\qquad d_i-\sigma_i d_{i-1}=0\quad(1\leq i\leq N).
\end{equation}
At node \(N\), attach a complete binary copy tree with
\begin{equation}
 M=2^{\lceil\log_2(8(N+1))\rceil},\quad K=2M-2,
 \quad P=N+1+K,
\end{equation}
and impose \(d_j-d_{\pi(j)}=0\) on every new node.  Thus all \(K\) copy
nodes have difference \(p_N\), while row and column sparsity remain at most
four and three.

We use the following elementary decoder throughout the section.

\begin{lemma}[Fixed-observable parity decoder]
\label{lem:lp-decoder}
Let \(O\) be an input-independent Hermitian observable with \(\norm O_2\leq1\).
If a family of normalized ideal states obeys
\[
 p_N\langle\psi_\sigma|O|\psi_\sigma\rangle\geq\delta>0,
\]
then the POVM \((I\pm O)/2\) guesses \(p_N\) with bias at least
\(\delta/2\).  From any \(\rho_\sigma\) with
\(\Dtr(\rho_\sigma,\ket{\psi_\sigma}\!\bra{\psi_\sigma})\leq\epsilon\),
the bias is at least \((\delta-2\epsilon)/2\).  If this is positive by an
absolute constant, preparing \(\rho_\sigma\) requires \(\Omega(N)\) raw
coefficient queries.
\end{lemma}

\begin{proof}
Trace duality changes the expectation of \(O\) by at most \(2\epsilon\).
The stated POVM has signed expectation \(\operatorname{Tr}(O\rho)\).  Constant
repetition turns any constant positive bias into bounded error.  Every sparse
row, column, position, or value query in the families below exposes at most one
\(\sigma_i\), even coherently, and is simulated by one call to the sign oracle.
The parity lower bound completes the reduction.
\end{proof}

First minimize \(\sum_jq_j\) subject to the chain and copy equations.  Exact
feasibility fixes every \(d_j\in\{\pm1\}\), and \(q_j\geq1\).  Hence the
unique optimum has \(q_j=1\) and value \(P\).  The nullspace has the public
orthonormal basis
\begin{equation}
 V_j=\frac{e_{u_j}+e_{v_j}}{\sqrt2}.
\end{equation}
On the feasible affine space the barrier objective is
\begin{equation}
 \Phi_\mu(q)=\sum_j\left[q_j-
       \mu\log\frac{q_j^2-1}{4}\right].
\end{equation}
Thus every central coordinate is
\(q(\mu)=\mu+\sqrt{\mu^2+1}\), and
\begin{equation}
 V^T\nabla^2\Phi_\mu V
 =2\mu\bigl[(q+1)^{-2}+(q-1)^{-2}\bigr]I.
\end{equation}
At \(\mu=3/4\), \(q=2\), so every oriented pair is
\((3/2,1/2)\).  If instead an exactly feasible point has objective gap at
most one, writing \(q_j=1+\delta_j\) gives
\(\sum_j\delta_j\leq1\) and
\begin{equation}
 \norm x_2^2=\sum_j\frac{q_j^2+1}{2}\leq P+\frac32.
\end{equation}
Since \(K\geq16(N+1)-2\), more than \(9/10\) of this mass lies in the
copy tree, and within an oriented copy pair the correct coordinate has at
least \(9/10\) of the pair mass.  The computational-basis decoder therefore
succeeds with probability greater than \(43/50\); trace error \(1/10\)
leaves success greater than \(19/25\).  This proves an \(\Omega(N)\) lower
bound both for every such near-optimal primal state and for the exact
\(\mu=3/4\) central state, despite \(\kappa_{\rm red}=1\).

There is a stronger any-feasible-point contract.  Add \(t_j\geq0\) and
\begin{equation}
 2u_j+2v_j+2t_j=3
\end{equation}
at every node, and retain objective \(\sum_jq_j\).  Exact feasibility now
gives
\begin{equation}
 1\leq q_j\leq\frac32,\qquad t_j=\frac32-q_j.
\end{equation}
The equality matrix has rank \(2P\), and its public orthonormal null basis is
\begin{equation}
 \bar V_j=\frac{e_{u_j}+e_{v_j}-2e_{t_j}}{\sqrt6}.
\end{equation}
The positive optimal columns form a basis; the opposite pair coordinate has
dual slack two.  Thus the LP is bounded and strictly primal--dual feasible,
with a unique nondegenerate strictly complementary optimum.  Its reduced
central Hessian is again a scalar identity, because the reduced barrier is a
sum of identical functions of \(q_j\).

For \(q\in[1,3/2]\), one triple has squared norm
\begin{equation}
 \bar w(q)=\frac{q^2+1}{2}+\left(\frac32-q\right)^2
 \in\left[\frac54,\frac{13}{8}\right].
\end{equation}
The copy triples therefore carry more than \(9/10\) of the total mass, and
the plus/minus decoder, randomized on \(t\), succeeds conditionally with
probability at least \(9/10\).  Consequently \emph{every} exact feasible
amplitude state, without an objective promise, has the same
\(\Omega(N)\) query lower bound at trace error \(1/10\).  The proof also
survives \(x\geq0\) and the absolute global \(\ell_1\) residual
\(\norm{\bar Ax-\bar b}_1\leq1/50\): telescoping gives
\(p_Nd_j\geq49/50\) at every copy node, while the cap gives
\(q_j+t_j\leq151/100\).  The resulting copy mass is greater than \(43/50\),
and the signed conditional expectation exceeds \(63/100\).  Before trace
error the total signed expectation is therefore
\begin{equation}
 \frac{43}{50}\frac{63}{100}=\frac{2709}{5000},
\end{equation}
so trace error \(1/10\) leaves signed expectation at least
\[
 \frac{2709}{5000}-\frac15=\frac{1709}{5000}>\frac13.
\]

\subsection{The robustness frontier}
\label{sec:lp-robustness-frontier}

The preceding \(\ell_1\) theorem cannot be promoted to a constant
\(\ell_2\) theorem on one unit-height path.  In gauged variables set
\begin{equation}
 p_id_i=1-\frac{i}{N},\qquad d_j=0\quad(j\text{ in the copy tree}),
\end{equation}
and take \(q_j=1,t_j=1/2\).  All cap and copy equations are exact, the
\(N\) chain residuals equal \(1/N\), and the total residual is
\(N^{-1/2}\), yet the amplified coordinates have no parity signal.  Using
\(1-2i/N\) instead gives an entirely wrong copy plateau at residual
\(2/\sqrt N\).  This is the basic drift obstruction.

A bounded-height repair, for power-of-two \(N\), uses \(R=N\) length-\(N\)
signed paths in parallel.
A binary fanout tree feeds the paths, and a binary tree with \(L=16N\)
leaves is attached to every endpoint.  The resulting tree has
\begin{equation}
 P=(2N-1)+N^2+N(2L-2)=33N^2-1,
 \qquad |S|=16N^2
\end{equation}
vertices and designated output leaves \(S\).  At every vertex use
\((u,v,h,t)\geq0\), impose pinned labelled-copy equations on
\(d=u-v\), pinned unsigned-copy equations on \(h\), and
\(q+t-2h=0\).  Minimize \(\sum_j(2u_j+2v_j+h_j+t_j)\).

\begin{theorem}[Parallel-path residual tube]
\label{thm:lp-parallel-path}
This LP has \(4P\) variables, \(3P\) full-rank equalities, coefficients of
magnitude at most two, and row and column sparsity at most four.  It has the
same regularity properties as the capped family and a scalar reduced central
Hessian for every \(\mu>0\).  If \(x\geq0\), \(x\ne0\), and
\begin{equation}
 \frac{\norm{Ax-b}_2}{\norm b_2}\leq\frac1{200},
 \qquad \norm b_2=\sqrt2,
\end{equation}
then the output-coordinate decoder has ideal bias greater than \(1/80\).
Consequently, preparing an amplitude state within trace distance \(1/100\)
of any such point needs \(\Omega(N)=\Omega(\sqrt P)\) raw queries.  Any
extra two-sided objective or centrality promise is immaterial.
\end{theorem}

\begin{proof}
Let \(\lambda_j\) be the exact propagated tree label.  Gauging removes all signs.
Descendant counting in each binary tree and Cauchy--Schwarz on the \(N\)
parallel paths give
\begin{equation}
 \frac{\norm{\operatorname{Diag}(\lambda)z-\mathbf1}_2}{\sqrt P}
 \leq12\norm{(z_r-1,Dz)}_2,
\quad
 \frac{\norm{\operatorname{Diag}(\lambda_S)z_S-\mathbf1}_2}{\sqrt{|S|}}
 \leq12\norm{(z_r-1,Dz)}_2.
\label{eq:lp-tag-9-15}
\end{equation}
The binary-level operator norms sum to
\(\sum_{\ell\geq1}2^{-\ell/2}=1+\sqrt2<5/2\), while the path term is
\(\sqrt{N/R}=1\).  With \(E=\norm{Ax-b}_2\leq\sqrt2/200\),
\eqref{eq:lp-tag-9-15} for
\(d\) and \(h\), plus the cap and nonnegativity, gives
\begin{equation}
 \norm x_2^2<11P,
 \qquad
 \sum_{j\in S}\lambda_jd_jq_j>\frac35|S|.
\end{equation}
The two loss terms in the second inequality, divided by \(|S|\), are at most
\[
 12E\left(1+2.18\sqrt{P/|S|}\right)
 <\frac{12\sqrt2}{200}\left(1+2.18\sqrt{33/16}\right)<\frac25.
\]
Hence the measurement bias is larger than
\[
 \frac{(3/5)16}{2\cdot11\cdot33}=\frac8{605}>\frac1{80}.
\]
Trace error \(1/100\) leaves positive constant bias, and
Lemma~\ref{lem:lp-decoder} applies.  Full rank and regularity follow by
pivoting successively on the pinned \(d\), pinned \(h\), and private
\(t\) blocks.  The public null vectors are the local vectors in
\eqref{eq:lp-gain-null}; identical feasible intervals make the central
curvature scalar.
\end{proof}

The quadratic volume is necessary in the signed-edge gadget class, not in all
bounded-coefficient real-linear gadgets.

\begin{proposition}[Effective-resistance volume bound]
\label{prop:lp-resistance-volume}
Let a connected graph have maximum degree \(\Delta\), a pinned root, and
outputs \(S\).  Suppose every consistent edge constraint
\(d_v-a_{uv}(\sigma)d_u=0\) has \(a_{uv}\in\{\pm1\}\) depending on at most
\(k\) input bits, and every output label is parity.  Then there is a vector
with root value one and wrong-parity value at every output whose residual is
at most
\begin{equation}
 \frac{k\sqrt{2\Delta|V|}}{N}.
\label{eq:lp-tag-9-17}
\end{equation}
Thus excluding all wrong-output vectors below an absolute residual \(\eta\)
requires \(|V|>\eta^2N^2/(2\Delta k^2)\).
\end{proposition}

\begin{proof}
The product of edge labels on a root-to-output path is parity.  Fourier degree
is subadditive under products, so every such path has length at least \(N/k\).
Gauge the consistent labels away and consider a unit root-to-\(S\) flow.
Removing circulations and decomposing into paths gives
\(\sum_e|f_e|\geq N/k\).  Cauchy--Schwarz and Thomson's principle yield
\[
 R_{\rm eff}(r,S)\geq\frac{N^2}{k^2|E|},\qquad
 C_{\rm eff}(r,S)\leq\frac{\Delta k^2|V|}{2N^2}.
\]
The Dirichlet energy for voltage difference two is \(4C_{\rm eff}\); its
harmonic minimizer lies in \([-1,1]\).  Taking the square root and undoing
the gauge proves \eqref{eq:lp-tag-9-17}.  The next section develops the general mixture and
cut frontiers; only this signed-edge specialization is needed here.
\end{proof}

There is a load-bearing escape witness.  On one path replace the signed
recurrence by
\(d_i-2\sigma_i d_{i-1}=0\), add a public reference
\(h_i-2h_{i-1}=0\), and retain \(q_i+t_i-2h_i=0\).  The exact scale is
\(h_i=|d_i|=2^i\).  With absolute residual
\(E=\norm{Ax-b}_2\), rescaling by \(2^{-i}\) gives
\begin{equation}
 \left|2^{-i}p_id_i-1\right|
 \leq\left(1+\frac1{\sqrt3}\right)E.
\end{equation}
Put \(K_0=1+1/\sqrt3\), and assume \(E\leq10^{-2}\).  Applying the same
estimate to the reference rows gives
\[
 p_id_i\geq(1-K_0E)2^i,\qquad
 0\leq h_i\leq(1+K_0E)2^i.
\]
If \(c_i=q_i+t_i-2h_i\) is the cap residual, nonnegativity gives
\(q_i,t_i\leq2h_i+|c_i|\).  Therefore
\begin{equation}
 \norm x_2^2
 \leq\sum_i\bigl[(q_i+t_i)^2+h_i^2\bigr]
 \leq12(1+K_0E)^2\,4^N+2E^2,
\qquad
 x_{{\rm correct},N}\geq(1-K_0E)2^N.
\end{equation}
The endpoint consequently has constant normalized mass (and ideal decoder
bias greater than \(3/100\) at \(E=10^{-2}\)), giving a linear-size robust
family.  It has exponential range and poor conditioning, showing that bounded
local coefficients alone do not imply the quadratic volume law.  The next
construction replaces the long gain chain by a short gain prefix and a
polynomial plateau.

\subsection{The linear-size gain--plateau family}
\label{sec:lp-gain-plateau}

Set
\begin{equation}
 T=\left\lceil\frac12\log_2N\right\rceil,
 \quad K=16N,\quad P=17N+1,
 \quad H_i=2^{\min\{i,T\}},\quad H=2^T\in[\sqrt N,2\sqrt N).
\end{equation}
There are nodes \(i=0,\ldots,N+K\), gains
\(g_i=H_i/H_{i-1}\in\{1,2\}\), and labels \(a_i=\sigma_i\) for
\(i\leq N\) and \(a_i=1\) afterward.  Define
\[
 \tau_0=1,\qquad \tau_i=\prod_{j=1}^i a_j,\qquad
 R_\tau=\diag(\tau_0,\ldots,\tau_{P-1}).
\]
At every node introduce
\((u_i,v_i,h_i,t_i)\geq0\), put \(d_i=u_i-v_i\), \(q_i=u_i+v_i\), and
impose
\begin{equation}
\begin{aligned}
 d_0&=1,&d_i-g_ia_id_{i-1}&=0,\\
 h_0&=1,&h_i-g_ih_{i-1}&=0,\\
 &&q_i+t_i-2h_i&=0.
\end{aligned}
\label{eq:lp-tag-9-20}
\end{equation}
The base objective is \(\sum_i(2u_i+2v_i+h_i+t_i)\).  Its public
orthonormal null basis is
\begin{equation}
 W_i=\sqrt{\frac23}\left(\frac12e_{u_i}+\frac12e_{v_i}-e_{t_i}\right).
 \label{eq:lp-gain-null}
\end{equation}

\begin{theorem}[Linear robust primal-state lower bound]
\label{thm:lp-gain-plateau}
The LP \eqref{eq:lp-tag-9-20} has \(4P\) variables, \(3P\) full-rank equalities,
coefficients of magnitude at most two, and row and column sparsity at most
four and three.  It is bounded and strictly primal--dual feasible, with a
unique nondegenerate strictly complementary optimum.  The optimum is
\(\Theta(N^{3/2})\), its largest coordinate is \(\Theta(\sqrt N)\), and
its Euclidean norm is \(\Theta(N)\).

For every nonzero \(x\geq0\) with
\begin{equation}
 \operatorname{res}_p(x)\leq\frac1{100},
\label{eq:lp-tag-9-21}
\end{equation}
the fixed plateau decoder has ideal bias greater than \(1/30\).  Preparing
its original-primal state to trace error \(1/100\) therefore needs
\(\Omega(N)=\Omega(P)\) raw queries.  Moreover,
\begin{equation}
 \kappa_{\rm red}(\mu)<\frac76\qquad(0<\mu\leq1/16).
\label{eq:lp-tag-9-22}
\end{equation}
\end{theorem}

\begin{proof}
The pinned signed and unsigned bidiagonal blocks fix
\(d_i=H_i\tau_i\) and \(h_i=H_i\); private \(t_i\) columns finish the rank
proof.  Exact feasibility is
\begin{equation}
 H_i\leq q_i\leq2H_i,\qquad t_i=2H_i-q_i.
\end{equation}
The local objective is \(q_i+3H_i\), so the unique optimum has
\(q_i=t_i=H_i\).  The sign-selected pair member together with \(h_i,t_i\)
forms a triangular positive basis, and the opposite coordinate has slack two.
This proves all regularity claims.  Since there are \(\Theta(N)\) nodes of
height \(H=\Theta(\sqrt N)\), the three scale estimates follow.

Write a central coordinate as \(q_i=H_i+z_i\).  Stationarity and curvature
are
\begin{equation}
 \frac1\mu=\frac1{2H_i+z_i}+\frac1{z_i}-\frac1{H_i-z_i},
\quad
 f_i''=\mu\left[(2H_i+z_i)^{-2}+z_i^{-2}+(H_i-z_i)^{-2}\right].
\end{equation}
For \(\mu\leq1/16\),
\(15\mu/16\leq z_i<\mu\), whence
\(1/\mu<f_i''<7/(6\mu)\).  Changing to the orthonormal coordinates
\eqref{eq:lp-gain-null} multiplies every eigenvalue by \(2/3\), proving
\eqref{eq:lp-tag-9-22}.

For robustness put \(E=\norm{Ax-b}_2\leq\sqrt2/100\).  The weighted
resistance is
\[
 \sum_{i\geq1}H_i^{-2}<\frac13+\frac{17N}{H^2}\leq\frac{52}{3}.
\]
Weighted telescoping gives, simultaneously at every node,
\begin{equation}
 \left|\frac{\tau_id_i}{H_i}-1\right|<\frac3{40},
 \qquad \left|\frac{h_i}{H_i}-1\right|<\frac3{40}.
\end{equation}
The cap then gives \(q_i,t_i<2.165H_i\), and
\begin{equation}
 \norm x_2^2<11\sum_iH_i^2
 <\frac{583}{3}NH^2.
\end{equation}
On the \(K=16N\) output nodes,
\(p_Nd_iq_i>(37H/40)^2\).  The computational-basis decoder has bias
\begin{equation}
 \frac{16N(37/40)^2H^2}{2(583/3)NH^2}>\frac1{30}.
\end{equation}
Trace error \(1/100\) leaves bias greater than \(7/300\), and
Lemma~\ref{lem:lp-decoder} completes the proof.
\end{proof}

The early-path conditioning claim cannot be strengthened: for \(\mu/H\to
\infty\), local curvature scales as \(\mu/H_i^2\), so the unrigidified
condition ratio can reach \(\Theta(H^2)=\Theta(N)\).  The gain is also the
precise escape from the quadratic path-bundle law.  If \(R\) internally
disjoint length-\(N\) paths have reference height at most \(H\), a harmonic
wrong-endpoint interpolation has squared residual at most \(4RH^2/N\).
Therefore exclusion below absolute residual \(\eta\) requires
\begin{equation}
 RH^2>\frac{\eta^2N}{4},\qquad
 P=\Theta(RN)\Longrightarrow PH^2=\Omega(N^2).
\label{eq:lp-tag-9-28}
\end{equation}
Both the bounded-height parallel construction and the linear-volume
gain--plateau construction attain this frontier.  The general gadget-class
mixture and cut theorems are given in \cref{sec:gadget-frontiers};
\eqref{eq:lp-tag-9-28} is
the only
path-bundle statement used here.  Naive balanced-XOR/PCP encodings do not
improve it: boundary truth-table points obstruct strict feasibility, and
homogenized shallow circuits still admit fractional drift across a small cut.

For one static input, Theorem~\ref{thm:lp-gain-plateau} gives the tight total
ledger
\begin{equation}
 Q_{\rm prep}+\sum_{t=1}^R
 (Q_{{\rm form},t}+Q_{{\rm solve},t}+Q_{{\rm recover},t})+Q_{\rm out}
 =\Omega(P).
\label{eq:lp-tag-9-29}
\end{equation}
The bound is total, not per iteration.  Querying all \(N\) signs and caching
their prefixes costs \(O(P)\) and makes every later raw query free, so no
\(\Omega(RP)\) fixed-instance theorem is possible.

\subsection{Reduced geometry versus the affine offset}
\label{sec:lp-affine-separation}

All reduced geometry in the gain--plateau family is public.  Every exact
feasible point has the form
\begin{equation}
 x_\sigma(q)_i=
 \left(\frac{q_i+H_i\tau_i}{2},
       \frac{q_i-H_i\tau_i}{2},H_i,2H_i-q_i\right),
\label{eq:lp-tag-9-30}
\end{equation}
and
\begin{equation}
 \{x:A_\sigma x=b\}=x_\sigma^{\rm off}+\operatorname{range}(W),
\quad
 (x_\sigma^{\rm off})_i
 =H_i\left(\frac{1+\tau_i}{2},\frac{1-\tau_i}{2},1,1\right).
\end{equation}
Even the minimum-norm offset is hidden: minimizing the local squared norm in
\eqref{eq:lp-tag-9-30} gives \(A_\sigma^\dagger b=x_\sigma(4H_i/3)\).  By contrast, the
full reduced barrier, all derivatives, central scalars, \(W\), both
projectors, and the diagonal reduced Hessian are independent of \(\sigma\).
For \(\mu\leq1/16\), its eigenvalues lie in
\((2/(3\mu^2),7/(9\mu^2))\) in the unscaled log-barrier convention of
\cref{def:reduced-kappa}; the \(\mu\)-weighted Hessian
\(W^\top\nabla^2\Phi_\mu W\) has eigenvalues in
\((2/(3\mu),7/(9\mu))\).

\begin{corollary}[Affine-offset oracle dichotomy]
\label{cor:lp-affine-dichotomy}
Grant arbitrary exact input-independent reduced-geometry oracles for free.
For the gain--plateau family, loading any original-primal state satisfying
\eqref{eq:lp-tag-9-21} costs \(\Omega(P)\), while the free reduced Hessian has condition
below \(7/6\) on the late tail.  For the parallel-path family, the analogous
cost is \(\Omega(\sqrt P)\), while the reduced Hessian is an exact scalar
identity for the entire central path.  In both cases the hidden object is the
affine offset, not the reduced solution.
\end{corollary}

\begin{proof}
Every granted channel is the same for all \(\sigma\) and can be hardwired in
the parity reduction.  Theorem~\ref{thm:lp-gain-plateau} and
Theorem~\ref{thm:lp-parallel-path} then apply unchanged.  A free feasible
offset, prefix table, or loading isometry would not be input-independent: one
known output pair reveals \(p_N\), and a normalized offset state carries a
constant amount of output-pair mass.
\end{proof}

\subsection{Exact Newton secants}
\label{sec:lp-newton-secants}

The base objective in the gain--plateau family gives a uniformly conditioned
late central tail, but its exact public-start construction initially used
different complementarity parameters on the \(O(\log N)\) gain nodes.  A
public reweighting removes that qualification.  At height \(H_i\), replace
the local objective by
\begin{equation}
 (1+\omega_i)(u_i+v_i)+h_i+t_i,
 \qquad \omega_i=\frac7{36H_i}.
\end{equation}
On the feasible line its nonconstant part is \(\omega_iq_i\), so the unique
optimum and all regularity properties survive; the nonbasic reduced cost is
\(2\omega_i>0\).

Set
\begin{equation}
 \mu_1=\frac1{16},\qquad
 x_i^1=H_i\left(\frac98,\frac18,1,\frac34\right)
\label{eq:lp-tag-9-33}
\end{equation}
up to the hidden pair swap, and put \(s_j^1=\mu_1/x_j^1\).  The reduced
stationarity equation at \(z_i=H_i/4\) is exact because
\begin{equation}
 \frac{\omega_i}{\mu_1}=\frac{28}{9H_i}
 =\frac1{2H_i+H_i/4}+\frac1{H_i/4}-\frac1{H_i-H_i/4}.
\end{equation}
Thus \((x^1,y^1,s^1)\) is the genuine global \(\mu_1\)-central point.

Use the public start
\begin{equation}
\begin{split}
 x_i^0&=H_i(5/6,5/6,20/27,5/9),\\
 s_i^0&=H_i^{-1}(10/27,10/27,5/12,5/9),\qquad y^0=0.
\end{split}
\label{eq:lp-tag-9-35}
\end{equation}
Every coordinate product is
\begin{equation}
 \mu_0=\frac{25}{81},\qquad
 \widehat\mu=\frac{25}{162}=\frac{\mu_0}{2}.
\end{equation}
Every Newton right-hand side is public, because the pair differences at the
start vanish.  For either hidden pair orientation,
\begin{equation}
 \frac{10}{27H_i}\frac{9H_i}{8}
 +\frac{5H_i}{6}\frac1{18H_i}
 =\frac{10}{27H_i}\frac{H_i}{8}
 +\frac{5H_i}{6}\frac1{2H_i}
 =\frac{25}{54}.
\label{eq:lp-tag-9-37}
\end{equation}
The \(h,t\) coordinates give the same \(25/54\).  Hence
\begin{equation}
 S^0(x^1-x^0)+X^0(s^1-s^0)
 =\left(\frac{25}{54}-2\frac{25}{81}\right)e
 =(\widehat\mu-\mu_0)e.
\label{eq:lp-tag-9-38}
\end{equation}
Primal and dual endpoint feasibility provide the other two rows of the
standard system \eqref{eq:full-kkt}; full row rank makes the correction
unique.  One full algebraic Newton step therefore reaches the genuine global
\(\mu_1=1/16\) center.  The endpoint parameter differs from
\(\widehat\mu\) because the Newton equation omits
\(\Delta X\Delta s\).

Up to the hidden swap, the primal direction block is
\begin{equation}
 \Delta x_i=H_i(7/24,-17/24,7/27,7/36),
 \qquad \norm{\Delta x_i}_2^2
 =D H_i^2,\quad D=\frac{16139}{23328}.
\label{eq:lp-tag-9-39}
\end{equation}
The unselected-minus-selected squared pair gap is \(5H_i^2/12\) (the
unselected coordinate carries the larger direction mass, \(17H_i/24\)
against \(7H_i/24\)).  Since
\(\sum_iH_i^2<(17N+4/3)H^2\), the fixed output-copy decoder has bias
\begin{equation}
 \frac{16N(5/12)H^2}{2D(17N+4/3)H^2}>\frac14.
\label{eq:lp-tag-9-40}
\end{equation}
Thus the exact original-primal direction state at trace error \(1/100\)
needs \(\Omega(P)\) raw queries.

In the unscaled log-barrier convention of
\cref{def:reduced-kappa}, the \(\mu_1\)-scaled endpoint reduced Hessian and
the start reduced Newton matrix are
\begin{equation}
 \mu_1 W^T\nabla^2F(x^1)W
 =\frac{182}{243}\diag(H_i^{-2}),
\qquad
 W^T(X^0)^{-1}S^0W=\frac{22}{27}\diag(H_i^{-2}).
\end{equation}
Thus all but the \(T=O(\log N)\) public gain modes form an exact scalar
block.  The public congruence \(D_H=\diag(H_i)\) makes both matrices scalar;
as a standard-form diagonal change of variables it raises the coefficient
scale to \(\Theta(\sqrt N)\).  More generally, scaling by \(H_i^\eta\)
gives
\begin{equation}
 B_\eta=\Theta(H^\eta),\quad
 \kappa_\eta=H^{2(1-\eta)},\quad
 B_\eta^2\kappa_\eta=\Theta(N).
\end{equation}
This is a property of this diagonal family, not a universal
normalization--condition theorem.

The remaining outliers are forced inside this local ansatz.

\begin{proposition}[Four-variable-node incompatibility]
\label{prop:lp-local-incompatibility}
Consider hidden-swappable endpoint pairs with values \(L_i>S_i>0\) at one
common central parameter \(\mu_1\), and suppose their differences are the
prescribed heights \(L_i-S_i=H_i\).  Within the separable four-variable
node gadget, the following three requirements are incompatible: a public
start centered coordinatewise at one \(\mu_0\); one full standard Newton
correction with a common target that works for both hidden orientations; and
an \(O(1)\)-conditioned endpoint Hessian in the raw orthonormal null basis.
This statement does not cover added rows, weighted barriers, or other local
gadgets.
\end{proposition}

\begin{proof}
Write the two public start values as \(a_i,b_i\), with slacks
\(\mu_0/a_i,\mu_0/b_i\).  Equality of the first-coordinate secant cross
terms for the two hidden orientations requires
\[
 \frac{\mu_0L_i}{a_i}+\frac{a_i\mu_1}{L_i}
 =\frac{\mu_0S_i}{a_i}+\frac{a_i\mu_1}{S_i}.
\]
Since \(L_i\ne S_i\), this gives
\(a_i^2=\mu_0L_iS_i/\mu_1\); the second coordinate similarly gives
\(b_i=a_i\).  The common cross term is therefore
\[
 C_i=\sqrt{\mu_0\mu_1}\,
       \frac{L_i+S_i}{\sqrt{L_iS_i}}.
\]
One global Newton target makes \(C_i\) independent of \(i\), hence the
aspect ratio \(L_i/S_i\) is constant.  The cap row
\(t_i=2H_i-(L_i+S_i)>0\) makes \(L_i+S_i\) proportional to \(H_i\) and
forces that constant aspect ratio to exceed three.  Together with
\(L_i-S_i=H_i\), this makes both endpoint values proportional to \(H_i\).
The log-barrier curvature in the local orthonormal null direction
consequently scales as \(H_i^{-2}\), giving condition ratio
\(H^2=\Theta(N)\).
\end{proof}

The reweighted construction also exposes an important negative result.  A
fixed relative residual for the full Newton system alone does not force the
primal direction ray.  On the \(K\) output nodes let
\(w=\sum_{i\in S}W_i\), and perturb the exact solution by
\begin{equation}
 e_x=Mw,\qquad e_y=0,\qquad
 e_s=-(X^0)^{-1}S^0e_x.
\end{equation}
The primal and complementarity residual perturbations vanish.  On the
plateau, direct substitution in \eqref{eq:lp-tag-9-35} gives the sharp bound
\[
 \norm{(X^0)^{-1}S^0w}_2\leq\frac{\norm w_2}{H^2}
\]
(the exact coefficient is below \(0.856/H^2\)), whereas the public cap RHS
has norm \(\Omega(H\sqrt K)\).  With \(M=H^2\), the relative full-KKT
residual is \(O(H^{-1})\) and the normalized primal direction converges to
the public zero-query state \(\ket{w/\norm w}\).  The perturbed update has
negative \(t\)-coordinates.  Therefore the robust direction theorem below
must retain the geometric condition \(x^0+\widetilde{\Delta x}\geq0\); a
bare fixed full-KKT residual promise is insufficient.

\subsubsection{Prefix rigidification: the condition-one capstone}
\label{sec:lp-prefix-rigidified}

Return to the constraints \eqref{eq:lp-tag-9-20}.  At each of the \(T\) small-height nodes
\(i<T\), add the public row
\begin{equation}
 q_i-\frac54h_i=0,
\label{eq:lp-tag-9-44}
\end{equation}
and use the uniform objective
\begin{equation}
 \sum_i[(1+\omega)(u_i+v_i)+h_i+t_i],
 \qquad \omega=\frac7{36H}.
\end{equation}
Write the augmented matrix and RHS as \(A'_\sigma,b'\).

\begin{theorem}[Condition-one original-primal Newton hardness]
\label{thm:lp-prefix-capstone}
The prefix-rigidified family has \(4P\) variables and \(3P+T\) full-rank
equalities.  Coefficients have magnitude at most two, and row and column
sparsity are at most four.  It is bounded and strictly primal--dual feasible,
with a unique nondegenerate strictly complementary optimum.  In the raw
orthonormal reduced coordinates,
\begin{equation}
 \kappa_{\rm red}(\mu)=1\quad\text{for every }\mu>0,
\end{equation}
and the reduced Newton matrix at the public start \eqref{eq:lp-tag-9-35} is also exactly
condition one.

The unique standard Newton correction from that start with target
\(\widehat\mu=25/162\) lands after one full step at the genuine global
\(\mu_1=1/16\) center \eqref{eq:lp-tag-9-33}.  Its original-primal direction state has
decoder bias greater than \(1/4\), so trace error \(1/100\) requires
\(\Omega(N)=\Omega(P)\) raw queries.

The same lower bound holds for a state of every nonzero \(x\geq0\) satisfying
\(\operatorname{res}_p(x)\leq1/100\), and for every nonzero direction
\(\widetilde{\Delta x}\) whose update is nonnegative and obeys that residual,
at state trace error \(1/100\).  In the latter case a fixed interference
decoder has bias greater than \(1/54\).
\end{theorem}

\begin{proof}
For the row \(r_i\) in \eqref{eq:lp-tag-9-44} and the old null vectors,
\begin{equation}
 r_i^TW_j=\sqrt{2/3}\,\delta_{ij}.
\end{equation}
Hence the new rows are independent modulo the old row space and
\begin{equation}
 \ker A'_\sigma=\operatorname{span}\{W_i:i\geq T\}.
\label{eq:lp-tag-9-48}
\end{equation}
They fix \(q_i=5H_i/4\) on the prefix, leaving only equal-height \(H\)
coordinates free.  At the optimum all four prefix variables and three
variables at every free node are positive, exactly \(3P+T\) in total.  A
null vector supported on them must vanish because every nonzero free \(W_i\)
touches that node's zero pair coordinate.  The positive columns form a basis,
and each nonbasic reduced cost is \(2\omega>0\).  This proves rank and
regularity.

All surviving nodes have the same height, objective slope, central scalar,
and local curvature.  Equation~\eqref{eq:lp-tag-9-48} therefore gives
\[
 W^T\nabla^2F(x(\mu))W=\lambda(\mu)I_{P-T}.
\]
At \(\mu_1=1/16\), the free coordinate is \(q=5H/4\), while the fixed
prefix coordinates already have the same local template.  Positive slacks
\(s^1=\mu_1(X^1)^{-1}e\) are orthogonal to the surviving kernel after
subtracting the objective, so full row rank supplies the central multiplier.
The secant identities \eqref{eq:lp-tag-9-37}--\eqref{eq:lp-tag-9-38} are unchanged, and the endpoint satisfies
the new rows.  This proves the exact Newton statement.  At the start,
\begin{equation}
 W^T(X^0)^{-1}S^0W=\frac{22}{27H^2}I_{P-T}.
\end{equation}
The fraction-to-boundary maximum is exactly
\(\alpha_{\max}=20/17>1\), so the full step preserves strict positivity;
any safety factor at least \(17/20\) permits it.  The direction formula
\eqref{eq:lp-tag-9-39} survives and \eqref{eq:lp-tag-9-40} proves the exact-state lower bound.

For a robust feasible state, deleting the public new residual coordinates
gives the base residual bound with the same RHS norm
\(\norm{b'}_2=\sqrt2\); Theorem~\ref{thm:lp-gain-plateau} applies.  For a
robust direction, set \(\widetilde x=x^0+\widetilde{\Delta x}\geq0\).  On
every plateau node the same telescoping estimates give
\begin{equation}
 \tau_i\widetilde d_i>\frac{37}{40}H,\quad
 \frac{37}{40}H<\widetilde h_i<\frac{43}{40}H,\quad
 \widetilde q_i,\widetilde t_i<\frac{433}{200}H.
\label{eq:lp-tag-9-50}
\end{equation}
Together with the public start, these imply
\begin{equation}
 \norm{\widetilde{\Delta x}}_2^2<100NH^2,\qquad
 \widetilde{\Delta h}_i>\frac{199}{1080}H
 \quad(i\in S).
\label{eq:lp-tag-9-51}
\end{equation}
Let \(\ket{g_i}=(\ket{u_i}-\ket{v_i})/\sqrt2\), and use the public norm-one
observable
\begin{equation}
 O=\bigoplus_{i\in S}
 (\ket{h_i}\!\bra{g_i}+\ket{g_i}\!\bra{h_i}).
\end{equation}
Equations \eqref{eq:lp-tag-9-50}--\eqref{eq:lp-tag-9-51} give signed expectation greater than
\begin{equation}
 \frac{7363\sqrt2}{270000}>\frac1{27},
\end{equation}
so its associated binary POVM has bias greater than \(1/54\).  Trace error
\(1/100\) leaves positive constant bias.  Lemma~\ref{lem:lp-decoder}
finishes the proof.
\end{proof}

The condition-one claim is intentionally reduced-primal.  At the same start,
the raw OSS matrix \({\cal M}_0=[-X^0(A')^T\ \ S^0W]\) and the eliminated
augmented KKT matrix satisfy
\begin{equation}
 \kappa_2({\cal M}_0)\geq\sqrt{681/44}\,H^2,
 \qquad
 \kappa_2({\cal K}_0)\geq\sqrt{1701/178}\,H^2,
\end{equation}
so both are \(\Omega(P)\).  Exact block preconditioning can reduce the
second condition to \(\varphi^2\), but its difference-row Schur block has
the gauge form
\begin{equation}
 G_\sigma^{(d)}=R_\tau G_+^{(d)}R_\tau,\qquad
 (G_\sigma^{(d)})^{-1/2}
 =R_\tau(G_+^{(d)})^{-1/2}R_\tau.
\end{equation}
The all-positive block is an irreducible tridiagonal Stieltjes matrix, whose
inverse square root is entrywise strictly positive by its resolvent integral.
Thus the sign of a sign-resolving preconditioner entry reveals a prefix
product.  A factor block encoding without entry access requires a separate
implementation analysis; \cref{sec:preconditioning} treats that
factor-access frontier.

The access localization can be exact.  In the Moore--Penrose gauge write
\begin{equation}
 \Delta x=p_\sigma+Wz,\qquad
 p_\sigma=(A'_\sigma)^\dagger(b'-A'_\sigma x^0).
\end{equation}
On every free node,
\begin{equation}
 (p_d,p_q,p_h,p_t)=H(\tau_i,-4/27,7/27,-2/27),
\end{equation}
with pair coordinates recovered from \(p_d,p_q\).  The complete reduced
system is public:
\begin{equation}
 M=\frac{22}{27H^2}I,\qquad
 z_i=-\frac{29}{108}\sqrt{\frac32}\,H,\qquad Mz=g.
\end{equation}
Normalized at \(\alpha_M=22/(27H^2)\), it has
\(M/\alpha_M=I\), \(\kappa=1\), and
\begin{equation}
 \alpha_M\frac{\norm{M^{-1}g}_2}{\norm g_2}=1.
\end{equation}
The states of \(g,z,Wz\) are public and free.  For the fixed diagonal-pair
observable
\[
 O_S=\sum_{i\in S}
 \bigl(\ket{u_i}\!\bra{u_i}-\ket{v_i}\!\bra{v_i}\bigr),
\]
each output node contributes
\(u_i^2-v_i^2=p_d p_q=-(4/27)p_NH^2\).  A free-node block has squared norm
\((851/1458)H^2\), a prefix block has squared norm
\(DH_i^2\) with \(D=16139/23328<7/10\), and
\(\sum_iH_i^2<(53/3)NH^2\).  Hence
\begin{equation}
 \left|\left\langle
 \frac{p_\sigma}{\norm{p_\sigma}_2},O_S
 \frac{p_\sigma}{\norm{p_\sigma}_2}
 \right\rangle\right|
 >\frac{(64/27)NH^2}{(7/10)(53/3)NH^2}
 =\frac{1920}{10017}>\frac16.
\end{equation}
Its sign is \(-p_N\), so
preparing \(\ket{p_\sigma/\norm{p_\sigma}}\) to trace error \(1/100\)
costs \(\Theta(P)\) raw queries.  Because
\(p_\sigma\perp Wz\), their coherent combination succeeds with probability
at least \(1/2\), provided one has phase-referenced coherent preparation
access to both normalized summands.  A density-state contract alone does not
supply that access, as \cref{sec:prelim-output-contracts} records.  The
localization is gauge-specific: an arbitrary
particular solution can trade a null component with \(z\); the invariant
claim is the end-to-end original-coordinate lower bound.

The capstone has the same tight static ledger \eqref{eq:lp-tag-9-29}.  Its matching upper
bound caches all signs in \(N\) queries.  A genuine iteration multiplier is
valid only for a streaming service in which round \(t\) releases a fresh
independent sign block after output \(t-1\) is committed.  Before the next
block is released, round \(t\) must return an unconditional state within
trace distance \(1/100\) of that round's hard original-coordinate Newton
direction (or of a point under the robust contract in
\cref{thm:lp-prefix-capstone}).  Let \(Q_t\) count every call to the current
block's coefficient oracle or inverse, all setup after that block is
released, and all queries used to prepare classical or quantum advice
consumed in that round.  Conditioned on all past memory, the new block is
independent, so \cref{thm:lp-prefix-capstone} applies in every round and
\begin{equation}
 \sum_{t=1}^RQ_t=\Omega(RN)=\Omega(RP).
\end{equation}
This is a freshness direct sum, not a per-iteration theorem for one fixed LP.

\subsection{Full-KKT and preconditioned formulations}
\label{sec:lp-full-kkt}

The preceding results request only the original-primal block.  A related
construction makes the complete primal--dual direction hard while keeping its
unpreconditioned right-hand side public.  Use the unrigidified gain--plateau
constraints \eqref{eq:lp-tag-9-20}, the base objective \(c_i=(2,2,1,1)\), and the public start
\begin{equation}
 x^0=s^0=e,\qquad y^0=0,\qquad \mu^0=1.
\end{equation}
For a public centering target \(\theta\in[0,1]\), the Newton equations are
\begin{equation}
 A_\sigma\Delta x=b-A_\sigma e,\quad
 A_\sigma^T\Delta y+\Delta s=c-e,\quad
 \Delta x+\Delta s=(\theta-1)e.
\label{eq:lp-tag-9-62}
\end{equation}
The RHS is input-independent: every hidden coefficient multiplies the zero
pair difference at the start.

Let \(R_{r,j}=\sum_{i=j}^{N+K}H_i^r\).  Direct elimination and backward
substitution give
\begin{equation}
\begin{aligned}
 \Delta x_{u_i}&=(4H_i-8+3H_i\tau_i)/6,&
 \Delta x_{v_i}&=(4H_i-8-3H_i\tau_i)/6,\\
 \Delta x_{h_i}&=H_i-1,&
 \Delta x_{t_i}&=(2H_i-1)/3,\\
 \alpha_j&=\tau_jR_{2,j}/(2H_j),&
 \beta_j&=[7R_{2,j}+(4-9\theta)R_{1,j}]/(3H_j),\\
 \gamma_j&=(2H_j+2)/3-\theta .
\end{aligned}
\label{eq:lp-tag-9-63}
\end{equation}
Here \((\alpha,\beta,\gamma)=\Delta y\) correspond to the signed,
reference, and cap rows.

\begin{theorem}[Public-RHS full-KKT direction]
\label{thm:lp-full-kkt}
Let \(w_\sigma=(\Delta x,\Delta y,\Delta s)\) solve \eqref{eq:lp-tag-9-62}.
Preparing its
normalized state to trace distance \(1/100\) takes at least \(N/2\) raw
coefficient queries.  The same holds if an approximate vector obeys
\(\norm{\widetilde w-w_\sigma}_2\leq\norm{w_\sigma}_2/300\) and its state
has trace error at most \(1/300\).

For \(\theta=1\), the two-block elimination has an explicit canonical sparse
block encoding of normalization six and a public RHS state.  The same
\(N/2\) lower bound holds for calls to that fully specified block encoding,
not merely to its principal matrix block.
\end{theorem}

\begin{proof}
On the \(K\) output rows define
\(A_S=H^2\sum_{\ell=1}^K\ell^2/4\).  The height sums and
\eqref{eq:lp-tag-9-63} imply
\begin{equation}
 \norm\alpha_2^2<3A_S,\qquad
 \norm\beta_2^2<\frac{484}{3}A_S,\qquad
 \norm{w_\sigma}_2^2<165A_S.
\end{equation}
Moreover \(\alpha_j\) has sign \(p_N\) and \(\beta_j>0\) on the output
suffix, with \(\beta_j/|\alpha_j|\geq3\).  The public norm-one observable
that swaps the corresponding \(\alpha_j,\beta_j\) coordinates therefore has
parity-signed expectation
\begin{equation}
 p_N\langle O_{\rm KKT}\rangle>\frac6{165}=\frac2{55}>\frac1{30}.
\label{eq:lp-tag-9-65}
\end{equation}
Trace error \(1/100\) changes this by at most \(1/50\).  The resulting
expectation is a real polynomial of degree at most twice the query count; its
nonzero parity Fourier coefficient forces degree at least \(N\).  Thus the
query count is at least \(N/2\).  The triangle inequality for normalized
rays gives the approximate-vector statement.

At \(\theta=1\), eliminate \(\Delta s=-\Delta x\), put \(y'=-\Delta y\),
and write
\begin{equation}
 {\cal K}_\sigma=\begin{bmatrix}I&A_\sigma^T\\A_\sigma&0\end{bmatrix},
 \qquad
 r_\sigma=\binom{e-c}{b-A_\sigma e}.
\end{equation}
The maximum absolute row and column sums are six.  The standard
magnitude-weighted row/column state-pair construction gives an exact normalization-six
block encoding.  Its support, magnitudes, failure arms, and left preparation
are public; the right preparation uses one local sign phase, hence exactly
one sign query per call, adjoint, or controlled call.  Also
\(\norm{r_\sigma}_2^2=3P+T+1\), and its preparation is public.  Removing
\(\Delta s\) can only strengthen \eqref{eq:lp-tag-9-65}, so the same degree argument applies
to this complete oracle implementation.
\end{proof}

Ideal Schur preconditioning makes the matrix well conditioned but does not
make the end-to-end raw-input task easy.  Put
\[
 {\cal S}_\sigma=A_\sigma A_\sigma^T,\quad
 Q_\sigma={\cal S}_\sigma^{-1/2}A_\sigma,\quad
 P_\sigma=\diag(I,{\cal S}_\sigma^{-1/2}).
\]
Define
\[
 \overline w_\sigma=P_\sigma^{-1}\binom{\Delta x}{y'},\qquad
 \overline r_\sigma=P_\sigma r_\sigma,\qquad
 \overline K_\sigma=P_\sigma{\cal K}_\sigma P_\sigma.
\]
Then
\begin{equation}
 \overline K_\sigma=
 \begin{bmatrix}I&Q_\sigma^T\\Q_\sigma&0\end{bmatrix},
 \qquad
 \overline K_\sigma\overline w_\sigma=\overline r_\sigma,
\end{equation}
and the primal block of \(\overline w_\sigma\) is exactly \(\Delta x\).
Since \(Q_\sigma Q_\sigma^T=I\), the spectrum is \(1\) on the nullspace and
\begin{equation}
 \frac{1+\sqrt5}{2},\qquad \frac{1-\sqrt5}{2}
\end{equation}
on every coupled two-plane.  Therefore
\begin{equation}
 \norm{\overline K_\sigma}_2
 =\norm{\overline K_\sigma^{-1}}_2=\varphi,\qquad
 \kappa_{\rm abs}(\overline K_\sigma)=\varphi^2
 =\frac{3+\sqrt5}{2}.
\label{eq:lp-tag-9-69}
\end{equation}
This is the classical ideal-Schur spectrum of Murphy--Golub--Wathen
\cite{murphy2000-a-note-on-preconditioning-for}.

\begin{theorem}[Residual-robust Schur-preconditioned KKT state]
\label{thm:lp-preconditioned-kkt}
Suppose \(\widetilde{\overline w}\ne0\) obeys
\begin{equation}
 \norm{\overline K_\sigma\widetilde{\overline w}
       -\overline r_\sigma}_2
 \leq10^{-3}\norm{\overline r_\sigma}_2.
\label{eq:lp-tag-9-70}
\end{equation}
Preparing its normalized state to trace error \(1/200\) takes at least
\(N/2\) total raw coefficient queries, counting preconditioner, transformed
RHS, solve, and recovery access.  The result transfers to the
prefix-rigidified family.
\end{theorem}

\begin{proof}
Let
\[
 O_x=\sum_{i\in S}
 \bigl(\ket{D_i}\!\bra{h_i}+\ket{h_i}\!\bra{D_i}\bigr),
 \qquad
 \ket{D_i}=\frac{\ket{u_i}-\ket{v_i}}{\sqrt2}.
\]
This norm-one, primal-supported observable has exact parity-signed
expectation greater than \(2/87\) on the preconditioned solution.  Indeed,
\begin{equation}
 \norm{\Delta x}_2^2\leq\frac{11}{3}PH^2,\qquad
 p_N\frac{\langle\Delta x|O_x|\Delta x\rangle}
              {\norm{\Delta x}_2^2}>\frac16,\qquad
 \norm{\overline w}_2^2<\frac{29}{4}\norm{\Delta x}_2^2.
\end{equation}
Equations \eqref{eq:lp-tag-9-69}--\eqref{eq:lp-tag-9-70} give
\(\norm{\widetilde{\overline w}-\overline w}_2
 \leq10^{-3}\varphi^2\norm{\overline w}_2\).  The total target-state
distance is at most
\begin{equation}
 \frac1{200}+\frac{2\varphi^2}{1000}<\frac1{87},
\label{eq:lp-tag-9-72}
\end{equation}
which preserves the sign of the observable.  The polynomial argument gives
\(N/2\).

For the prefix-rigidified transfer, use the same all-ones start.  Every
added row \(q_i-\frac54h_i=0\) has the public primal residual \(-3/4\), and
the uniform objective is
\[
 c'_i=(1+\omega,1+\omega,1,1),\qquad
 \omega=\frac7{36H}.
\]
The hidden difference coefficients still multiply \(u_i-v_i=0\), so the
complete Newton RHS remains public.  On every free node, null projection and
exact feasibility give
\[
 \Delta q_i=\frac{4H-6-2\omega}{3},\qquad
 \Delta h_i=H-1,\qquad
 \Delta u_i-\Delta v_i=H\tau_i.
\]
The prefix variables are fixed by feasibility.  These formulas yield
\[
 \norm{\Delta x}_2^2<5\sum_iH_i^2,\qquad
 p_N\frac{\langle\Delta x|O_x|\Delta x\rangle}
              {\norm{\Delta x}_2^2}>\frac18,
\]
because the numerator is \(K\sqrt2H(H-1)\), each prefix node contributes
less than \(5H_i^2\), and each free node less than \(4H^2\).
For the exact block preconditioner the transformed primal RHS is
\(e-c'=(-\omega,-\omega,0,0)\) per node, and
\[
 \frac{\norm{e-c'}_2}{\norm{\Delta x}_2}<\frac{21}{100},
 \qquad
 \norm{\overline w_\sigma}_2^2
 <\frac52\norm{\Delta x}_2^2.
\]
Thus the parity-signed expectation on the complete preconditioned state
exceeds \(1/20\).  Its condition remains \(\varphi^2\), while the solve and
output errors in \eqref{eq:lp-tag-9-72} change the expectation by less than
\(2/87<1/20\).  The same \(N/2\) conclusion follows.
\end{proof}

There is no matrix-only normalization/condition tradeoff behind this theorem.
The row-space projector is public because \(\ker A_\sigma\) is public.  If
\(Q_0\) is a public row isometry with \(Q_0^TQ_0=I-WW^T\), then
\(O_\sigma=Q_\sigma Q_0^T\) is orthogonal and
\begin{equation}
 \diag(I,O_\sigma)^T\overline K_\sigma\diag(I,O_\sigma)
 =\begin{bmatrix}I&Q_0^T\\Q_0&0\end{bmatrix},
\end{equation}
which is wholly public.  Input dependence has moved to the transformed RHS
and recovery map.  Accordingly the invariant accounting statement is
\begin{equation}
 Q_{\rm pre}+Q_{\rm on}\geq N/2.
\end{equation}
It is tight: after \(N\) cached sign queries, later raw-query cost can vanish.
The separate factor-access frontier is given in
\cref{sec:preconditioning}; it is not proved here.

\subsection{Bounded data and the sharp residual scale}
\label{sec:lp-bounded-data}

Remove the gain prefix by setting every height and gain to one, retain a
length-\(N\) signed chain followed by \(K=16N\) public copy edges, and put
\(P=17N+1\).  Use the same four nonnegative variables at every node and the
equations
\begin{equation}
\begin{aligned}
 d_0&=1,&d_i-a_id_{i-1}&=0,\\
 h_0&=1,&h_i-h_{i-1}&=0,\\
 &&q_i+t_i-2h_i&=0,
\end{aligned}
\qquad
 a_i=\begin{cases}\sigma_i&i\leq N,\\1&i>N.\end{cases}
\end{equation}
Use objective
\begin{equation}
 c^Tx=\sum_i\left[\frac{43}{36}(u_i+v_i)+h_i+t_i\right].
\end{equation}
All matrix, RHS, objective, feasible-primal, and relevant slack coordinates
are bounded by absolute constants; \(\norm b_2=\sqrt2\).  Exact feasibility
is \(d_i=\tau_i,h_i=1,t_i=2-q_i,1\leq q_i\leq2\), so the same basis
\eqref{eq:lp-gain-null} spans the nullspace.

\begin{theorem}[Bounded exact-state separation]
\label{thm:lp-bounded-exact}
The all-unit family is bounded and strictly primal--dual feasible, with a
unique nondegenerate strictly complementary optimum.  Its raw orthonormal
reduced central Hessian has condition one for every \(\mu>0\), and its
reduced Newton matrix at the public start below is \((22/27)I\).

At \(\mu_1=1/16\), one exact Newton correction from
\begin{equation}
 x_i^0=(5/6,5/6,20/27,5/9),\quad
 s_i^0=(10/27,10/27,5/12,5/9),\quad y^0=0
\label{eq:lp-tag-9-77}
\end{equation}
with \(\mu_0=25/81\) and target \(25/162\) lands at the exact center
\begin{equation}
 x_i^1=(9/8,1/8,1,3/4),\qquad
 s_i^1=(1/18,1/2,1/16,1/12),
\end{equation}
up to the hidden pair swap.  Preparing either the exact primal direction
state or exact central-primal state to trace error \(1/100\) requires
\(\Omega(N)=\Omega(P)\) raw queries.
\end{theorem}

\begin{proof}
On the feasible line the local objective is \(3+(7/36)q_i\), so the unique
optimum has \(q_i=1\).  The positive columns form a basis and the zero pair
coordinate has reduced cost \(7/18\).  All central coordinates solve the
same scalar equation
\[
 \frac{7}{36\mu}=\frac1{q+1}+\frac1{q-1}-\frac1{2-q},
\]
so the reduced Hessian is scalar.  At \(q=5/4,\mu=1/16\), both sides are
\(28/9\), and direct evaluation gives the \(\mu_1\)-scaled reduced
curvature \(182/243\).  Equations
\eqref{eq:lp-tag-9-37}--\eqref{eq:lp-tag-9-38} with \(H_i=1\) prove the same exact Newton secant and
\((22/27)I\) start matrix.

The primal direction is the constant block in \eqref{eq:lp-tag-9-39}.  Its plateau decoder
has bias
\[
 \frac{K(5/12)}{2DP}>\frac14.
\]
Each central block has squared norm \(91/32\) and oriented pair gap \(5/4\),
so the central-state decoder has bias
\[
 \frac{K(5/4)}{2(91/32)P}=\frac{20K}{91P}>\frac15.
\]
Lemma~\ref{lem:lp-decoder} completes both lower bounds.
\end{proof}

The all-unit family is robust, but only at the sharp vanishing scale.

\begin{theorem}[Sharp all-unit residual threshold]
\label{thm:lp-unit-threshold}
For every nonzero \(x\geq0\) satisfying the absolute residual bound
\begin{equation}
 \norm{A_\sigma x-b}_2\leq\frac1{16\sqrt N},
\label{eq:lp-tag-9-79}
\end{equation}
the normalized primal state has constant correct-parity bias.  Preparing it
to trace error \(1/100\) needs \(\Omega(N)=\Omega(P)\) raw queries.  In
relative form, \eqref{eq:lp-tag-9-79} is \(1/(16\sqrt{2N})\); in particular it contains
the \(1/(100\sqrt N)\) absolute residual tube.

The order is sharp: a nonnegative zero-signal plateau exists at absolute
residual \(1/\sqrt N\), and a nonnegative entirely wrong-parity plateau
exists at absolute residual \(2/\sqrt N\), even with exact optimal objective.
After public homogeneous rigidification, the wrong witness can also have
exact dual feasibility and exact coordinatewise complementarity.
\end{theorem}

\begin{proof}
Let \(E=\norm{Ax-b}_2\), \(z_i=\tau_id_i\), and
\(\delta=\sqrt P E\).  Telescoping the signed and reference recurrences gives
\begin{equation}
 |z_i-1|\leq\sqrt P E,\qquad |h_i-1|\leq\sqrt P E.
\end{equation}
The factor \(\sqrt P\) is exact by equality in Cauchy--Schwarz.  Under
\eqref{eq:lp-tag-9-79}, \(P/N\leq35/2\) gives \(\delta<4/15\), so on every output node
\(p_Nd_i>11/15\).  The cap and nonnegativity give
\(q_i,t_i<31/12\), \(h_i<19/15\), and
\(\norm x_2^2<15P\).  Thus the output decoder has bias
\begin{equation}
 \frac{K(11/15)^2}{30P}>\frac1{70}.
\end{equation}
Trace error \(1/100\) leaves bias at least
\(1/70-1/100=3/700>0\).

For the zero-signal witness take
\begin{equation}
 z_i=1-i/N\ (i\leq N),\qquad z_i=0\ (i>N),
\end{equation}
and for the wrong-parity witness replace this by
\begin{equation}
 z_i=1-2i/N\ (i\leq N),\qquad z_i=-1\ (i>N).
\end{equation}
Set \(d_i=\tau_i z_i\) and \(h_i=q_i=t_i=1\).  All reference and cap rows
are exact.  The \(N\) signed residuals are respectively \(1/N\) and
\(2/N\), proving the claimed norms.  Both points are nonnegative and have
the exact optimal objective.

For the stronger failure, add \(q_i-5h_i/4=0\) on all non-output nodes and
average the exact centers of the \(N\) inputs obtained by flipping one sign.
The average has the wrong common output template and the same
\(2/\sqrt N\) primal residual.  Every coordinate is positive.  Setting
\(\bar s=\mu_1\bar X^{-1}e\) gives exact complementarity.  The remaining
nullspace is supported on output nodes, where \(\bar x\) is itself an exact
central template; hence \(c-\bar s\perp\ker A\), which supplies an exact
dual multiplier.
\end{proof}

The recurrence-only cutoff is exactly \(E<1/\sqrt P\): taking
\(z_i=1-(i+1)/P\) makes all gauged residuals \(-1/P\), reaches zero at the
last node, and has norm \(1/\sqrt P\).  Inverting the robust scale over all
hidden lengths \(L\) with \(17L+1\leq P\) gives the worst-case accuracy law
\begin{equation}
 Q(P,\epsilon)=\Omega\bigl(\min\{P,\epsilon^{-2}\}\bigr)
\end{equation}
for absolute or RHS-relative residual, up to constants.  This is a
reparameterized family bound, not accumulating hardness at several accuracies
on one instance.

The public copy plateau is necessary for constant one-copy distinguishability
at unit height.  Flipping only \(\sigma_N\) swaps exactly \(K+1\) target
blocks.  For the central and direction states the normalized overlaps are
\begin{equation}
 1-\frac{32(K+1)}{91P},\qquad
 1-\frac{K+1}{DP},
\end{equation}
respectively.  If \(K=o(N)\), their trace distance is \(o(1)\) even for
opposite parities.  Hence the correct bounded-data resource is
\(K=\Theta(N)\) copies, not \(\Theta(\sqrt N)\) dynamic range.

\subsection{Dual homogeneity and the bounded full direction}
\label{sec:lp-dual-homogeneity}

There is one exact way to suppress the multiplier accumulation of a unit-height
chain.  Positive dual scaling is an elementary symmetry of the central and
Newton equations, consistent with the general scale-invariance viewpoint in
interior-point theory \cite{tuncel1998-primal-dual-symmetry-scale-invariance}.

\begin{lemma}[Dual homogeneity]
\label{lem:lp-dual-homogeneity}
In the standard Newton system, leave \((A,b,x^0)\) fixed and replace
\begin{equation}
 (c,y^0,s^0,\widehat\mu)
 \longmapsto
 (\lambda c,\lambda y^0,\lambda s^0,\lambda\widehat\mu),
 \qquad\lambda>0.
\end{equation}
Then \((\Delta x,\Delta y,\Delta s)\) is replaced by
\((\Delta x,\lambda\Delta y,\lambda\Delta s)\).  The primal central
trajectory is unchanged, with \(\mu\) reparameterized as
\(\mu'=\lambda\mu\).
\end{lemma}

\begin{proof}
The primal equation is unchanged, while substitution multiplies the dual and
complementarity equations by \(\lambda\).  Full row rank and positivity make
the Newton solution unique.  The same substitution in
\(Ax=b,A^Ty+s=c,Xs=\mu e\) proves the central-path statement.
\end{proof}

For the full-state theorem use the all-unit signed chain and copy plateau, but
pin each reference coordinate independently by \(h_i=1\).  This changes the
row representation, not the primal feasible set.  Starting from
\eqref{eq:lp-tag-9-77}, scale
the objective, initial slacks, and complementarity data by
\begin{equation}
 \lambda=\frac1P.
\label{eq:lp-tag-9-87}
\end{equation}
The exact primal correction remains the block \eqref{eq:lp-tag-9-39}.  If the first coordinate
is paired with the larger endpoint coordinate, the unscaled slack correction
is
\begin{equation}
 \Delta\bar s_i=(-17/54,7/54,-17/48,-17/36),
\label{eq:lp-tag-9-88}
\end{equation}
and, with difference, height-pin, and cap multipliers,
\begin{equation}
 \Delta y=\frac1P(\bar\alpha,\bar\beta,\bar\gamma),\quad
 \bar\alpha_j=\frac29\tau_j(P-j),\quad
 \bar\beta_j=\frac{133}{48},\quad
 \bar\gamma_j=\frac{11}{12}.
\label{eq:lp-tag-9-89}
\end{equation}

\begin{theorem}[Bounded full-KKT direction]
\label{thm:lp-dual-homogeneous-full}
Every coordinate of the complete \(11P\)-dimensional direction
\(w=(\Delta x,\Delta y,\Delta s)\) in
\eqref{eq:lp-tag-9-87}--\eqref{eq:lp-tag-9-89} is bounded by an
absolute constant, and \(\norm w_2=\Theta(\sqrt P)\).  Preparing its
normalized state to trace error \(1/100\) requires
\(\Omega(N)=\Omega(P)\) raw queries.

The explicit canonical sparse block encoding of either the three-block Newton
matrix or its two-block elimination has normalization five, a public RHS, and
the same linear call lower bound.  The reduced Newton condition remains one,
but the central endpoint is at \(\mu_1'=1/(16P)\), and both raw full-KKT
matrices have condition \(\Theta(P)\).
\end{theorem}

\begin{proof}
Equations \eqref{eq:lp-tag-9-39}, \eqref{eq:lp-tag-9-88}, and
\eqref{eq:lp-tag-9-89} give
\[
 \norm{\Delta x}_\infty\leq17/24,
 \quad \norm{\Delta y}_\infty\leq2/9,
 \quad \norm{\Delta s}_\infty\leq17/(36P).
\]
Writing \(D=16139/23328\), direct summation gives
\begin{equation}
 DP\leq\norm w_2^2<\left(D+\frac1{40}\right)P<\frac34P.
\end{equation}
Use the diagonal decoder convention \(+1\) on every \(v\) coordinate and
\(-1\) on every \(u\) coordinate.  The underlying \(u\)-minus-\(v\)
squared-mass difference on each output node is \(-p_N(5/12)\), so the
decoder expectation contributes \(+p_N(5/12)\) per node and has magnitude
\begin{equation}
 \frac{(5/12)K}{(3/4)P}=\frac{5K}{9P}
 \geq\frac{32}{63}>\frac12.
\end{equation}
Trace error preserves constant bias.  The polynomial method gives at least
\(N/2\) sign queries.

The maximum absolute row and column sums of both canonical matrices are five.
The same magnitude-weighted state-pair construction used in
Theorem~\ref{thm:lp-full-kkt} therefore gives normalization five and one
sign query per full oracle call.  The RHS remains public because the start
has zero pair difference.  This proves the canonical-oracle claim.

Dual homogeneity gives \(\mu_1'=1/(16P)\) and the reduced start matrix
\((22/(27P))I\).  If \(z\in\ker A\) is unit, applying either raw KKT matrix
to the corresponding nullspace test vector has norm at most \(1/P\), while
the largest singular value is bounded below by a constant.  Conversely, the
pinned signed incidence has smallest singular value \(\Theta(P^{-1})\);
bounded local changes reduce the remaining blocks to identities.  The paired
saddle eigenvalues then show that both conditions are \(\Theta(P)\).
\end{proof}

The price is real: objective coefficients, positive slacks, reduced costs, and
complementarity are \(\Theta(P^{-1})\).  The theorem is not a simultaneous
constant-scale and constant-full-condition result.  It also does not inherit
the sharp relative-residual theorem, because independent height pins make
\(\norm b_2=\Theta(\sqrt P)\).

\subsection{Why the inverse-scale price remains open}
\label{sec:lp-beyond-pairs}

The obstruction can be stated without signed pairs.  At any positive start,
eliminate \(\Delta s\) from the Newton equations and put
\[
 H=(X^0)^{-1}S^0,\qquad g=(X^0)^{-1}r_c-r_d.
\]
Then
\begin{equation}
 H\Delta x-A^T\Delta y=g,\qquad
 \boxed{r_p^T\Delta y=\Delta x^TH\Delta x-g^T\Delta x}.
\end{equation}
This exact work identity is representation-independent.  If \(g=0\),
\(H\succeq\rho I\), and \(\norm{r_p}_1\leq R\), it gives
\begin{equation}
 \norm{\Delta y}_\infty\geq
 \frac{\rho\norm{\Delta x}_2^2}{R}.
\end{equation}
A constant-scale correction with \(\Theta(P)\) squared norm behind a sparse
public source therefore forces an \(\Omega(P)\) multiplier, independently of
redundant rows or branching.

There is also a two-input version for constant-dimensional copy blocks.  Let a
flat output region \(S\) carry \(z_v\in\R^k\) and orthogonal copy equations
\(z_v-U_ez_u=0\).  After public transports \(G_v\), suppose stationarity is
\begin{equation}
 \operatorname{div}_U(f)_v=D_v\Delta z_v+g_v,
 \qquad D_v\succeq\rho I,
\label{eq:lp-tag-9-94}
\end{equation}
with the same numerical \(D_v,g_v\) for two opposite-parity inputs, and their
transported codewords differ by a common \(w\), \(\norm w_2\geq2a\).  If a
boundary row scale is at most \(B\), summing the difference of
\eqref{eq:lp-tag-9-94} over
\(S\) cancels every public source and internal force.  Taking the inner
product with \(w\) yields
\begin{equation}
 \max_{p\in\{+,-\}}\norm{\Delta y_{\partial S}^p}_{2,\infty}
 \geq\frac{\rho a|S|}{B|\partial S|}.
\end{equation}
Thus swaps, simplex permutations, and public flat rotations do not remove
dual accumulation behind a bounded interface.

The most tempting counterexample is a real quarter-turn
\(J=\left[\begin{smallmatrix}0&-1\\1&0\end{smallmatrix}\right]\), with
\(z_i=\sigma_iJz_{i-1}\) and closure \(z_0-Jz_N=e_1\).  For \(4\mid N\),
\begin{equation}
 z_0(p)=\frac12(1,p)^T,
\end{equation}
so all blocks are bounded and locally reveal parity after a public inverse
rotation.  Nevertheless the work identity gives the closure multiplier
\begin{equation}
 e_1^T\Delta y_{\rm close}=\sum_{i=0}^N\norm{z_i}_2^2
 =\frac{N+1}{2}.
\label{eq:lp-tag-9-97}
\end{equation}
After gauging, the constraint is a twisted cycle whose singular values are
\begin{equation}
 2\left|\sin\frac{2\pi\ell\pm\pi/2}{2(N+1)}\right|,
 \qquad \ell=0,\ldots,N,
\end{equation}
so its least singular value is
\(2\sin(\pi/[4(N+1)])=\Theta(N^{-1})\).  A positive/negative standard-form
LP realization preserves bounded primal and slack directions and even allows
a positive half-centering step, but \eqref{eq:lp-tag-9-97} survives exactly.  Circulations,
public antisymmetric sources, and dense public residuals fail for the same
work or two-parity cancellation reason unless they introduce linearly many
genuinely independent hard interfaces.

\begin{openproblem}[Bounded-scale full-KKT parity amplifier]
\label{open:lp-bounded-full-kkt}
Does a sparse standard-form LP family exist that simultaneously has:
\begin{enumerate}[label=(\roman*)]
 \item linear dimension and bounded row and column degree;
 \item bounded coefficients with one-bit-local input dependence;
 \item public, simply preparable Newton right-hand sides;
 \item primal and slack starts bounded above and below by positive constants;
 \item a coordinatewise bounded complete Newton direction; and
 \item a parity-hard primal component with constant mass in that complete
 direction state?
\end{enumerate}
The work identity and orthogonal-copy cut law rule out the natural signed,
permutation, rotation, and flat-copy constructions.  They do not cover
arbitrary nonorthogonal mixed-mode constraints, input-dependent sources whose
loading is separately charged, or output interfaces that quotient the equality
multipliers.
\end{openproblem}

The distinction closing the section is therefore exact.  The
prefix-rigidified construction proves condition-one \emph{reduced Newton
geometry} together with a linear original-coordinate loading/recovery cost.
The preconditioned full-KKT result proves a constant-condition residual theorem
only after charging construction of its data-dependent access.  The bounded
full direction obtained by dual homogeneity instead has canonical full-system
condition \(\Theta(P)\).  None is a constant-\(\kappa\) QLS lower bound with
free matrix and RHS preparation oracles; each is an end-to-end LP access and
output theorem.
